\begin{table}[H]
\centering
\small
\setlength{\tabcolsep}{6pt}
\renewcommand{\arraystretch}{1.10}

\begin{tabular}{ccccc}
\toprule
\textbf{Outer Fold} & \textbf{$\delta$} & \textbf{$B^*$} & \textbf{LR Macro-F1} & \textbf{XGBoost Macro-F1} \\
\midrule
0 & 0.02 & 10 & 1.000 & 1.000 \\
1 & 0.02 & 10 & 1.000 & 1.000 \\
2 & 0.02 & 10 & 1.000 & 1.000 \\
3 & 0.02 & 10 & 1.000 & 1.000 \\
4 & 0.02 & 10 & 1.000 & 1.000 \\
\midrule
\textbf{All folds} & \textbf{0.02} & \textbf{10} & \textbf{1.000} & \textbf{1.000} \\
\bottomrule
\end{tabular}

\caption{
Outer-test-isolated feature-budget sensitivity for the mixed disease-associated signal. In each
outer fold, aptamer ranking and budget selection are performed using
outer-development patients only, and the selected budget is evaluated once on
the untouched outer-test patients. The ranking is fixed across inner validation
splits; $B$ is selected conditional on that ranking. Under the analysis-specified,
non-preregistered tolerance margin $\delta=0.02$, all five folds select the same development budget
$B^*=10$, and the resulting 10-aptamer panels preserve the observed
full-panel patient macro-F1 for both Logistic Regression and XGBoost.
This result demonstrates feature redundancy of the observed mixed
disease-associated signal. It does not validate a reduced assay, a biological
signature, or unique causal biomarkers.
}
\label{tab:panel_budget}

\end{table}
